
%The measurements included in this combination consist of the most recent Tevatron (Run 2) and LHC (Run 1 and Run 2) results.
The combination uses the latest measurements from D0 and CDF at the Tevatron and ATLAS and LHCb at the LHC.
The CMS Collaboration has not yet measured $m_W$, though it has measured differential $W$-boson cross sections
on the path to the measurement~\cite{CMSW}.
Prior measurements from the Tevatron and the CERN Super Proton Synchroton are not included as they are expected
to have negligible impact.  
The hadron-collider measurements are combined with
the result from the Large Electron Positron collider (LEP)~\cite{ALEPH:2013dgf},
$m_W = 80.376 \pm 0.033$~GeV\footnote{We use the convention $c \equiv 1$ and work in a right-handed coordinate system
with the origin at the centre of the detector and the $z$-axis along the beam pipe.  Cylindrical coordinates $(r,\phi)$
are used in the transverse plane, where $\phi$ is the azimuthal angle around the z-axis.  The pseudorapidity is defined
in terms of the polar angle $\theta$ as $\eta = -\ln\tan(\theta/2)$.  Transverse momentum is defined as $\pT = p\sin\theta$.
}.

The kinematic observables used in $m_W$ measurements at hadron colliders are the momentum of the charged lepton from
the $W$-boson decay ($p^{\ell}_{\scriptsize\textrm{T}}$) and the recoil transverse momentum ($u_{\scriptsize\mathrm{T}}$)
balancing the transverse momentum of the $W$ boson (\pTW).  The recoil is measured by vectorially summing the momentum of
all objects interacting in the detector, except for the charged lepton.  The neutrino momentum is inferred from the
net momentum imbalance,
$\vec{p}^{\nu}_{\scriptsize\mathrm{T}} \equiv -(\vec{p}^{\ell}_{\scriptsize\mathrm{T}} + \vec{u}_{\scriptsize\mathrm{T}})$.
For experiments with sufficiently good recoil resolution the most sensitive kinematic distribution is the transverse mass,
$m_{\scriptsize\textrm{T}} = \sqrt{2p_{\scriptsize\textrm{T}}^{\ell} p_{\scriptsize\textrm{T}}^\nu (1-\cos\Delta\phi)}$,
where $\Delta \phi$ is the angle between the charged lepton and the neutrino in the transverse plane.  

The CDF Collaboration measured $m_W$~\cite{CDF:2022hxs} using Run 2 data collected between 2003 and 2011 at the Tevatron
collider, corresponding to 8.8~\ifb{} of integrated luminosity from proton-antiproton ($p\bar{p}$) collisions at a
center-of-mass energy of $\sqrt{s}=1.96$~TeV.  The mass was obtained from template fits to the reconstructed
distributions of $p^{\ell}_{\scriptsize\mathrm{T}}$, \mT, and $p^{\nu}_{\mathrm T}$ in the 
electron and muon decay channels, yielding $m_W=80433.5\pm 6.4~\textrm{(stat.)} \pm 6.9~\textrm{(sys.)}$~MeV, or
$80433.5\pm 9.4$~MeV.  
The quoted value of $m_W$ corresponds to the NNPDF3.1 PDF set~\cite{Ball:2017nwa}, with the PDF uncertainty estimated using 
the largest 25 symmetric eigenvectors constructed through a principal-component analysis from the full replica
set.  The direct fit for $m_W$ to the data used events from a version of the {\textsc{ResBos}}~\cite{Ladinsky:1993zn}
generator referred to as {\textsc{ResBos-C}} in this paper.  The generation used the CTEQ6M PDF set~\cite{Pumplin:2002vw} 
and was tuned to fit the observed spectrum of $Z$-boson transverse momentum.  The uncertainty on the $W$-boson transverse 
momentum $p_{\scriptsize\textrm{T}}^{W}$ was determined using DYqT~\cite{Bozzi:2008bb,Bozzi:2010xn}, with a constraint from
the observed recoil
distribution in $W$-boson events.  The adjustment of the model to the NNPDF3.1 PDF set included an effective update of
the modelling of the leptonic angular distributions, as discussed in Sec.~\ref{sec:polarization}.

The D0 Collaboration performed two measurements of $m_W$ in Run 2 of the Tevatron collider.  The first used data taken
between 2002 and 2006, corresponding to an integrated luminosity of 1.1~\ifb{}~\cite{Abazov:2009cp}, and the second
used 2006-2008 data corresponding to an integrated luminosity of 4.3~\ifb{}~\cite{D0:2013jba}.  The analysis 
produced template fits for $m_W$ using the $p^{\ell}_{\scriptsize\textrm{T}}$, \mT, and $p^{\nu}_{\scriptsize\textrm{T}}$
kinematic distributions in the electron decay channel.  The initial 1.1~\ifb{} measurement combined the results from these
three distributions, while the 4.3~\ifb{} measurement removed the $p^{\nu}_{\scriptsize\textrm{T}}$ result due to its small
weight in the combination.  The overall combined result of all measurements is
$m_W = 80375\pm13~\textrm{(stat.)}\pm22~\textrm{(sys.)}$~MeV, or $80375\pm23$~MeV.  This value was determined using the
CTEQ6.1~\cite{Stump:2003yu} (CTEQ6.6~\cite{Nadolsky:2008zw}) PDF set for the measurement using 1.1~\ifb{} (4.3~\ifb{}).
The uncertainties were evaluated using {\textsc{Pythia6}}~\cite{Sjostrand:2006za} and the CTEQ6.1 PDF Hessian eigenvectors
scaled to reduce the nominal 90\% C.L. coverage to 68\% C.L.  The $p_{\scriptsize\textrm{T}}^{W}$ modelling used a version of
the {\textsc{ResBos}}~\cite{Balazs:1997xd,Landry:2002ix} generator referred to here as {\textsc{ResBos-CP}}.  

The $m_W$ measurement performed by the ATLAS Collaboration used $\sqrt{s}=7$~TeV proton-proton collision data
corresponding to 4.6~\ifb{} of integrated luminosity collected in 2011 during Run 1 of the LHC collider.
ATLAS performed template fits to the $p^{\ell}_{\scriptsize\textrm{T}}$ and \mT~distributions in the electron and muon
channels separately for
$W^+$ and $W^-$ events, since in proton-proton ($pp$) collisions the final-state distributions are different for these
processes.  The fits were further subdivided into three (four) pseudorapidity ranges in the electron (muon) channel,
yielding a total of 28 measurements.  The combination of these measurements yields 
$m_W=80370\pm7~\textrm{(stat.)}\pm18~\textrm{(sys.)}$, or $80370\pm19$~MeV.  The parton distribution functions were
modelled with the NNLO CT10 PDF set~\cite{Lai:2010vv}, with the Hessian uncertainties scaled to 68\% C.L.  The
$p_{\scriptsize\textrm{T}}^{W}$ modelling relied on the parton shower Monte Carlo (MC)
{\textsc{Pythia8}}~\cite{Sjostrand:2014zea} tuned to match the $p_{\scriptsize\textrm{T}}^{Z}$ distribution observed in
data.  The impact of the PDF uncertainties on the $m_W$ measurement was reduced by a simultaneous fit in different lepton
pseudorapidity regions.  The PDFs affect both the $p_{\scriptsize\textrm{T}}^{W}$ and $p_{\scriptsize\textrm{T}}^{Z}$
distributions, and to preserve the agreement with the $p_{\scriptsize\textrm{T}}^{Z}$ data distribution only the relative
variation between the $p_{\scriptsize\textrm{T}}^{W}$ and $p_{\scriptsize\textrm{T}}^{Z}$ distributions was propagated in
the uncertainty estimate.  Generated events were reweighted according to the calculation of the leptonic angular 
distributions in DYNNLO~\cite{Catani:2007vq,Catani:2009sm}.

The LHCb Collaboration performed a measurement of $m_W$ using Run 2 $pp$ LHC collision data collected in 2016 at 
$\sqrt{s}=13$~TeV, corresponding to 1.7~\ifb{} of integrated luminosity.  The measurement used the
$q/p^{\ell}_{\scriptsize\textrm{T}}$ distribution in the muon decay channels, where $q$ is the muon charge, giving a result
of $m_W = 80354 \pm 23~\textrm{(stat.)} \pm 10~\textrm{(exp.)} \pm 17~\textrm{(th.)}\pm 9~\textrm{(PDF)}$~MeV, or
$80354\pm32$~MeV.  The LHCb central value of $m_W$ and its uncertainty correspond to an unweighted average of results
using the
NNPDF3.1, MSHT2020~\cite{Bailey:2020ooq} and CT18~\cite{Hou:2019qau} PDF sets, all at next-to-leading order in the strong
coupling and with 68\% C.L. coverage.  The \pTW~distribution was modelled with
{\textsc{Powheg}~\cite{Nason:2004rx,Frixione:2007vw,Alioli:2010xd} interfaced to {\textsc{Pythia8}}, with a correction
at high boson $p_{\scriptsize\textrm{T}}$ derived from the observed $Z$-boson $p_{\scriptsize\textrm{T}}$ distribution.
The leptonic angular distributions were modelled with exact
$\mathcal{O}(\alpha_{\scriptsize\textrm{S}}^2)$ predictions from DYTurbo~\cite{Camarda:2019zyx}
and modified by scaling one of the leptonic angular coefficients when fitting the data.

The event requirements and fit ranges used in the measurements are summarized in Table~\ref{tab:evsel}.
CDF and D0 used similar analysis configurations, while at
ATLAS the looser recoil requirement and wider \mT~fit range were a consequence of the lower recoil resolution.
%The multijet background was enhanced in ATLAS as a consequence of this and the higher collision energy, though the
%narrower \pTl~fit range mitigated its impact.
The LHCb measurement was inclusive in recoil, with only a loose requirement on the momentum of the muon.
The ATLAS, CDF, and D0 measurements fit $m_W$ only, while LHCb performed a simultaneous fit for $m_W$ and the
relative fraction of $W^+$- to $W^-$-boson decays, the hadronic background fraction, $\alpha_{\scriptsize\textrm{S}}$
in $W$-boson events, $\alpha_{\scriptsize\textrm{S}}$ in $Z$-boson events, the intrinsic transverse momentum 
distribution of partons inside the proton, and the $A_3$ leptonic angular coefficient (see Sec.~\ref{sec:polarization}).

\begin{table}[!t]
  \centering
  \begin{tabular}{lll}
\hline
\hline
%    \toprule
    Experiment & Event requirements & Fit ranges \\
\hline
%    \midrule
    \vspace{-3mm}\\
    CDF   & $30<\pTl<55$~GeV & $32<\pTl<48$~GeV \\
          & $|\etal|<1$      & $32<\met<48$~GeV \\
          & $30<\met<55$~GeV & $60<\mT<100$~GeV \\
          & $65<\mT<90$~GeV  & \\
          & $\uT<15$~GeV  & \\
\vspace{-3mm}\\
    D0    & $\pT^e >25$~GeV  & $32<\pT^e<48$~GeV \\
          & $|\etal|<1.05$   & $65<\mT<90$~GeV  \\
          & $\met>25$~GeV    & \\
          & $\mT>50$~GeV     & \\ 
          & $\uT<15$~GeV     & \\
    \vspace{-3mm}\\
    ATLAS & $\pTl>30$~GeV  & $32<\pTl<45$~GeV \\
          & $|\etal|<2.4$  & $66<\mT<99$~GeV  \\
          & $\met>30$~GeV  & \\
          & $\mT>60$~GeV   & \\
          & $\uT<30$~GeV   & \\
    \vspace{-3mm}\\
    LHCb  & $\pT^\mu > 24$~GeV & $28 < \pT^\mu < 52$~GeV \\ 
          & $2.2 < \eta_\mu < 4.4$ & \\  
    \vspace{-3mm}\\
\hline \hline
%    \bottomrule
  \end{tabular}
  \caption{Event requirements and fit ranges for CDF, D0, ATLAS, and LHCb.\label{tab:evsel}}
\end{table}

